% =====================================================================
%  整理者　macanine <macanine@163.com>
%  仓库　　github.com/macanine/zju-linear-algebra-papers
%  来源　　解析据 papers/2025-2026-秋冬-期中-甲.tex 撰写（原卷见 sources/2025-2026-秋冬-期中-甲-回忆卷(含参考答案).pdf 第 1--2 页）
% =====================================================================

\begin{solution}
记系数矩阵的三个行向量为 $r_1,r_2,r_3$。未知数共 $4$ 个，通解含 $2$ 个任意常数，说明解空间的维数（即零空间的维数）为 $2$；由秩－零化度定理
\[
r+2=4\ \Longrightarrow\ r=2 .
\]
又 $\begin{vmatrix}1&1\\4&3\end{vmatrix}=3-4=-1\ne0$，故 $r_1,r_2$ 线性无关。秩为 $2$ 的三个行向量中 $r_3$ 必是 $r_1,r_2$ 的线性组合，设
\[
r_3=\lambda r_1+\mu r_2 .
\]
比较第 $2,3$ 个坐标：$\lambda+3\mu=1$，$\lambda+5\mu=3$，相减得 $2\mu=2$，故 $\mu=1$，再代回得 $\lambda=-2$。于是
\[
p=\lambda+4\mu=-2+4=2,\qquad -q=\lambda-\mu=-2-1=-3\ \Longrightarrow\ q=3 .
\]
此时第三式的右端由方程组自洽给出为 $\lambda\cdot(-1)+\mu\cdot(-1)=2-1=1$，恰等于题给右端 $1$，故有解且增广矩阵的秩也为 $2$。注意 $(p,q)$ 是被“秩为 $2$”与“有解”两条要求共同锁定的唯一取值。

取 $x_3,x_4$ 为自由未知量。前两个方程为
\[
x_1+x_2=-1-x_3-x_4,\qquad 4x_1+3x_2=-1-5x_3+x_4 .
\]
第一式乘 $4$ 后减去第二式得 $x_2=-3+x_3-5x_4$，再代回第一式得 $x_1=2-2x_3+4x_4$。
\ans{$p=2$，$q=3$，$r=2$，通解为
$X=\begin{pmatrix}2\\-3\\0\\0\end{pmatrix}+t_1\begin{pmatrix}-2\\1\\1\\0\end{pmatrix}+t_2\begin{pmatrix}4\\-5\\0\\1\end{pmatrix}$，$t_1,t_2\in\mathbb R$。}
\end{solution}

\begin{solution}
两矩阵相抵（等价）当且仅当秩相等，故先算两个行列式并定秩。按第一行展开得
\[
|A|=-(a-1)(a+1)=1-a^2,\qquad |B|=(a-1)(3a+2).
\]

\textbf{秩的分布。}$r(A)=3$ 当 $a\ne\pm1$；$a=1$ 时 $A$ 的第三行 $(4,4,4)$ 是第一行 $(1,1,1)$ 的 $4$ 倍，$a=-1$ 时第三行也是 $(4,4,4)$，故两种情形下都是 $r(A)\le2$；而 $A$ 的前两行 $(1,1,1)$、$(1,0,2)$ 线性无关（$\begin{vmatrix}1&1\\1&0\end{vmatrix}=-1\ne0$），故这两种情形下都是 $r(A)=2$。
$r(B)=3$ 当 $a\ne1,-\dfrac23$；$a=1$ 时 $B$ 的第三行 $(4,0,4)$ 是第一行 $(1,0,1)$ 的 $4$ 倍，$a=-\dfrac23$ 时 $|B|=0$，故这两种情形下都是 $r(B)\le2$；而 $B$ 的前两行 $(1,0,1)$、$(1,3,3)$ 线性无关（$\begin{vmatrix}1&0\\1&3\end{vmatrix}=3\ne0$），故这两种情形下都是 $r(B)=2$。逐值比较：
\begin{enumerate}[label=(\arabic*),leftmargin=2.2em,itemsep=0.25em,topsep=0.3em]
\item $a=1$：$r(A)=r(B)=2$；
\item $a=-1$：$r(A)=2\ne3=r(B)$；
\item $a=-\dfrac23$：$r(A)=3\ne2=r(B)$；
\item 其余 $a$：$r(A)=r(B)=3$。
\end{enumerate}
所以不相抵只可能发生在 (2)(3)，即 $a=-1$ 或 $a=-\dfrac23$；秩同为 $2$ 只在 (1)，即 $a=1$。

\textbf{解方程。}第 (4) 种情形下 $r(A)=3$，$|A|\ne0$，解唯一。直接代入 $X=(1,-1,1)^{T}$：
\[
A\begin{pmatrix}1\\-1\\1\end{pmatrix}
=\begin{pmatrix}1-1+1\\1+0+2\\4-4+a^2+3\end{pmatrix}
=\begin{pmatrix}1\\3\\a^2+3\end{pmatrix},
\]
故 $x_1=1$，$x_2=-1$，$x_3=1$。
\ans{$a=-1$ 或 $a=-\dfrac23$ 时 $A,B$ 不相抵；$a=1$ 时二者的秩都是 $2$；秩都是 $3$ 时 $x_1=1$，$x_2=-1$，$x_3=1$。}
\end{solution}

\begin{solution}
由展开定理，$D$ 按第一列展开为
\[
D=a_{11}A_{11}+a_{21}A_{21}+a_{31}A_{31}+a_{41}A_{41}.
\]
本题第一列的元素全为 $1$，即 $a_{11}=a_{21}=a_{31}=a_{41}=1$，而题设又给出 $A_{11}=A_{21}=A_{31}=A_{41}=1$，故直接得
\[
D=1\cdot1+1\cdot1+1\cdot1+1\cdot1=4 .
\]
一般地，$A_{i1}$ 与第一列的元素无关，所以 $\displaystyle\sum_{i=1}^{4}c_iA_{i1}$ 等于把 $D$ 的第一列换成 $(c_1,c_2,c_3,c_4)^{T}$ 后所得的行列式。本题取 $c_i=1$ 时换出的列恰是原来的全 $1$ 列，该行列式就是 $D$，两种算法给出同一结果。
\ans{$D=4$。}
\end{solution}

\begin{solution}
第一行元素 $a,b,c,d$ 可以是任意的，$A_{1j}$ 与它们无关，只由去掉第一行后剩下的 $3$ 阶子式决定，$x,z$ 也正是出现在这三行里。按代数余子式的定义逐项展开（$(1,j)$ 元带符号 $(-1)^{1+j}$）：
\begin{align*}
A_{11}&=-(x+3)(xz+x+2z+4), & A_{12}&=-(x+3)(xz+3x+2z+2),\\
A_{13}&=-xz-5x-2z, & A_{14}&=(x-1)(x+2)(x+3).
\end{align*}
同一组 $x,z$ 必须同时满足四个给定值。先看 $A_{14}$：由 $(x-1)(x+2)(x+3)=-4$，若 $x=-3$ 则左端为 $0$，矛盾，故 $x\ne-3$；于是 $A_{12}=0$ 可约去因子 $x+3$，化为
\[
xz+3x+2z+2=0 .
\]
又 $A_{13}=4$ 即 $xz+5x+2z=-4$。两式相减得 $2x=-2$，故 $x=-1$；代回 $xz+3x+2z+2=0$ 得 $-z-3+2z+2=0$，即 $z=1$。

回代检验：$A_{11}=-(2)(-1-1+2+4)=-8$，$A_{14}=(-2)(1)(2)=-4$，与题设相符。
\ans{$x=-1$，$z=1$。}
\rem{不要试图用 $D=\sum_j a_{1j}A_{1j}$ 去定 $x,z$——第一行是任意的，这条路走不通；正确的入口是“子式与第一行无关”以及同一组 $x,z$ 要让四个值同时成立。}
\end{solution}

\begin{solution}
$A^{*}$ 是上三角矩阵，其行列式等于对角元之积：
\[
|A^{*}|=1\cdot1\cdot1\cdot x=x .
\]
对 $n=4$ 阶矩阵有伴随矩阵的行列式公式 $|A^{*}|=|A|^{n-1}=|A|^{3}$，于是
\[
|A|^{3}=x .
\]
又 $n\ge2$ 时 $A$ 可逆当且仅当 $A^{*}$ 可逆：$A$ 可逆时 $A^{*}=|A|A^{-1}$ 可逆；反之若 $A$ 不可逆则 $|A|=0$，从而 $|A^{*}|=|A|^{n-1}=0$，$A^{*}$ 不可逆。故
\[
A\text{ 可逆}\iff |A^{*}|\ne0\iff x\ne0 .
\]
当 $x\ne0$ 时由 $|A|^{3}=x$ 得 $|A|=\sqrt[3]{x}$（实数的立方根唯一）。若 $x=0$，则 $|A|=0$，$A$ 不可逆，两种说法一致。
\ans{当 $x\ne0$ 时矩阵 $A$ 可逆；此时 $|A|=\sqrt[3]{x}$。}
\rem{$|A^{*}|$ 是 $x$ 的一次式，而 $|A^{*}|=|A|^{n-1}$ 中的指数是 $n-1=3$ 不是 $n$，把 $|A|$ 写成 $\sqrt[4]{x}$ 是常见错误。}
\end{solution}

\begin{solution}
$P$ 是初等矩阵，作用为“第 $3$ 行加上第 $1$ 行的 $2$ 倍”，即左乘 $P$ 实施 $R_3\leftarrow R_3+2R_1$；$Q$ 是交换第 $1,3$ 行（同时交换第 $1,3$ 列）的置换矩阵，满足
\[
Q^{2}=E\ \Longrightarrow\ Q^{2k+1}=Q,\quad Q^{2k}=E\qquad(k\ge0),
\]
即 $Q$ 的奇数次幂为 $Q$、偶数次幂为 $E$。$2025$ 为奇数、$1000$ 为偶数，故 $Q^{2025}=Q$，$Q^{1000}=E$。

先算 $PA$：只动第三行，
\[
PA=\begin{pmatrix}a_{11}&a_{12}&a_{13}\\a_{21}&a_{22}&a_{23}\\a_{31}+2a_{11}&a_{32}+2a_{12}&a_{33}+2a_{13}\end{pmatrix}.
\]
右乘 $Q^{2025}=Q$ 即交换第 $1,3$ 列，得
\[
PAQ^{2025}=\begin{pmatrix}a_{13}&a_{12}&a_{11}\\a_{23}&a_{22}&a_{21}\\a_{33}+2a_{13}&a_{32}+2a_{12}&a_{31}+2a_{11}\end{pmatrix}.
\]

再算 $P^{1025}$。记 $E_{31}$ 为仅 $(3,1)$ 元为 $1$、其余为 $0$ 的矩阵，则 $P=E+2E_{31}$，且 $E_{31}^{2}=O$，故二项式展开只剩前两项：
\[
P^{k}=(E+2E_{31})^{k}=E+2kE_{31}\qquad(k\ge1),
\]
即 $P^{k}$ 总是“第 $3$ 行加上第 $1$ 行的 $2k$ 倍”；取 $k=1025$ 得系数 $2050$。又 $Q^{1000}=E$，于是
\[
P^{1025}AQ^{1000}=P^{1025}A=\begin{pmatrix}a_{11}&a_{12}&a_{13}\\a_{21}&a_{22}&a_{23}\\a_{31}+2050a_{11}&a_{32}+2050a_{12}&a_{33}+2050a_{13}\end{pmatrix}.
\]
\ans{$PAQ^{2025}=\begin{pmatrix}a_{13}&a_{12}&a_{11}\\a_{23}&a_{22}&a_{21}\\a_{33}+2a_{13}&a_{32}+2a_{12}&a_{31}+2a_{11}\end{pmatrix}$，\\[0.2em]
$P^{1025}AQ^{1000}=\begin{pmatrix}a_{11}&a_{12}&a_{13}\\a_{21}&a_{22}&a_{23}\\a_{31}+2050a_{11}&a_{32}+2050a_{12}&a_{33}+2050a_{13}\end{pmatrix}$。}
\end{solution}

\begin{solution}
置换矩阵由 $E$ 经行、列互换得到，故它每行、每列恰有一个 $1$，其余元素为 $0$。于是 $A$ 的列向量是某个排列下的标准单位向量
\[
A=(e_{p_1}\ e_{p_2}\ \cdots\ e_{p_n}),
\]
其中 $p_1,\dots,p_n$ 是 $1,\dots,n$ 的一个排列。由 $(A^{T}A)_{ij}$ 等于 $A$ 的第 $i$ 列与第 $j$ 列的内积，
\[
(A^{T}A)_{ij}=e_{p_i}^{T}e_{p_j}=\delta_{ij}
=\begin{cases}1,&i=j,\\0,&i\ne j,\end{cases}
\]
故 $A^{T}A=E$。几何上这正说明 $A$ 的列是一组标准正交基，$A$ 是正交矩阵。
\ans{$A^{T}A=E$。}
\end{solution}

\begin{solution}
记 $\alpha=(a_1,a_2,\dots,a_n)^{T}$，$\beta=(b_1,b_2,\dots,b_n)^{T}$，则 $A=\alpha\beta^{T}$，题设即
\[
\beta^{T}\alpha=a_1b_1+a_2b_2+\cdots+a_nb_n=1 .
\]
由矩阵乘法的结合律
\[
A^{2}=(\alpha\beta^{T})(\alpha\beta^{T})=\alpha\bigl(\beta^{T}\alpha\bigr)\beta^{T}.
\]
其中 $\beta^{T}\alpha$ 是 $1\times1$ 的数，可以与向量交换次序提到最前面，故
\[
A^{2}=\bigl(\beta^{T}\alpha\bigr)\alpha\beta^{T}=1\cdot A=A .
\]
\qed
\end{solution}

\begin{solution}
由 $BC=0$，对每个 $j=1,\dots,n$ 都有
\[
B(Ce_j)=(BC)e_j=0,
\]
即 $C$ 的每一列都是齐次方程组 $BX=0$ 的解。因此 $C$ 的列空间含于 $\ker B$：
\[
\operatorname{Col}(C)\subseteq\ker B\ \Longrightarrow\ \dim\ker B\ge\dim\operatorname{Col}(C)=r(C).
\]

再由秩－零化度定理，$\dim\ker A=n-r(A)$。结合 $r(A)<r(C)$ 得
\[
\dim\ker A=n-r(A)>n-r(C)\ge n-\dim\ker B,
\]
即 $\dim\ker A+\dim\ker B>n$。$\ker A$ 与 $\ker B$ 都是 $\mathbb R^{n}$ 的子空间，由维数公式
\[
\dim(\ker A\cap\ker B)=\dim\ker A+\dim\ker B-\dim(\ker A+\ker B)
\ge\dim\ker A+\dim\ker B-n>0 .
\]
故 $\ker A\cap\ker B$ 中含非零向量 $X_0$，它满足 $AX_0=0$ 且 $BX_0=0$，从而 $AX_0=BX_0$。取 $n$ 阶矩阵 $X=(X_0,0,\dots,0)$（只有第一列为 $X_0$，其余列为零），则 $X\ne O$ 且
\[
AX=BX=0 .
\]
\qed
\rem{关键在于把“$BC=0$”翻译成“$C$ 的列都在 $\ker B$ 里”，于是 $\dim\ker B\ge r(C)>r(A)=n-\dim\ker A$；条件 $r(A)<r(C)$ 的作用正是把两个零空间的维数之和顶过 $n$，从而交集非零。}
\end{solution}
